% % Describe what the experiment(s) will look like. Depending on your study, it can include information like:
% %%% Research design
% %%% Dataset (including information on source, size, labelling, etc., to the extent available)
% %%% The different variations of the approach that you will compare, as well as details on how they will be applied (library, parameter tuning, etc.) 
% %%% Participants (numbers, intended characteristics)
% %%% Experimental procedure
% %%% Evaluation procedure / metrics

\section{Experimental Setup}

The research questions can be classified as evaluating the success of this setup or, more broadly, of typologically guided cross-lingual transfer (\textbf{Q1}, \textbf{Q2}, \textbf{Q3}) and as investigating how shifts in the model parameters facilitate cross-lingual transfer (\textbf{Q4}). Section~\ref{sub:evaluation} outlines how the former group is answered, Section~\ref{sub:probing} does the same for the latter. Table~\ref{tab:data} in Appendix~\ref{app:b} lists the specific datasets used in the experiments.

	\subsection{Model Evaluation}
	\label{sub:evaluation}
		Successful transfer is defined as higher performance on the target language test set. Following the example of Lauscher et al.~\cite{lauscher-etal-2020-zero}, the unlabeled attachment score (UAS) is used to quantify performance.\footnote{Using the more widely reported LAS would inhibit comparisons between datasets, as different treebanks may have been labeled with different annotation schemes.} In this section, $\text{UAS}($\mbase$, L)$ is taken to be the UAS of \mbase evaluated on the test set of language $L$. Cosine similarity between URIEL's \texttt{syntax\_knn} vectors for languages $L1$ and $L2$ is denoted as $\text{s}(L1, L2)$. Moreover, as stated previously, $\ell$ is used for any of the auxiliary languages, $\lambda$ is used for any of the target languages.

		\subsubsection{Q1}
		\textbf{Does greater typological similarity between the source and target languages result in better zero-shot cross-lingual transfer performance?}

			If typological similarity indeed leads to better zero-shot performance, the expectation is that $\text{UAS}($\mde$, l) - \text{UAS}($\mbase$, l)$ is higher when $\text{s}(\textit{deu}, l)$ is higher. That is, finetuning on German improves zero-shot performance more for typologically similar languages. Comparison with \mbase is key, since models might generally struggle with certain treebanks regardless of typology. Here, $l$ is any language in the union of auxiliary and target languages.

		\subsubsection{Q2}
		\textbf{To what degree does few-shot learning improve performance over zero-shot transfer?}
		
			While Lauscher et al.~\cite{lauscher-etal-2020-zero} found few-shot learning to yield considerable improvements over zero-shot performance, it is important to quantify the degree to which that holds in this setup. To eliminate presence in M-BERT's pre-training data as a factor, this question is only answered using few-shot learning on the target languages $\lambda$. 

			Following Stage 3, there are twenty models \m{\textit{\scriptsize l}}{\lambda} with $l$ now being either German or any of the auxiliary languages. For each, improvement induced by few-shot learning is defined as $\text{UAS}($\m{\textit{\scriptsize l}}{\lambda}$, \lambda) - \text{UAS}(M_{\textit{\scriptsize l}}, \lambda)$. This difference is expected to be positive. Differences between pairs of $l$ and $\lambda$ fall under the scope of \textbf{Q3}.

		\subsubsection{Q3}
		\textbf{Is few-shot transfer from a source to a target language more effective when first fine-tuning on a more typologically similar auxiliary language?}

			This question can effectively be split into two components. Both are answered by investigating whether $\text{UAS}($\m{l_1}{\lambda}$, \lambda) > \text{UAS}($\m{l_2}{\lambda}$, \lambda)$ whenever $\text{s}(l_1, \lambda) > \text{s}(l_2, \lambda)$. That is, few-shot performance is higher when the language finetuned on before few-shot learning on $\lambda$ is more typologically similar. One component of the question pertains to the effectiveness of the auxiliary stage in general. In this case, $l_1$ is any auxiliary language $\ell$ and $l_2$ is German. When $\ell$ is more similar to $\lambda$ than German is, the auxiliary stage may provide the model with new typological information that can then be leveraged during few-shot learning. The other component of the question pertains to the typological distance between the auxiliary language and the target language. In that case $l_1$ is any given auxiliary language and $l_2$ is any of the other auxiliary languages. This serves to assess whether any improvements found are the result of typology, rather than longer training regardless of typology.
		

	\subsection{Probing Experiments}
	\label{sub:probing}
		Following Choenni \& Shutova~\cite{choenni-shutova-2022-investigating}, investigation of the models' parameters is done using probing. In particular, they probe model embeddings for the presence of typological features. They first extract 10\,000 random sentences from the Tatoeba corpora for a number of different languages. Each sentence is passed through the language models they investigate, they use mean-pooling over the token embeddings to obtain a sentence embedding. The embedding is then passed through a one-hidden-layer MLP with 100 units, followed by an $n$-unit softmax output layer, where $n$ is the number of possible options for a given typological feature.

		For this thesis, a similar approach is used. However, since Tatoeba only has 404 sentences in Veps, 404 sentences are sampled for each of the nine languages used to keep the dataset balanced. The exact probing task is described below.

		\subsubsection{Q4}
		\textbf{Does fine-tuning on a language make its typological features more easily identifiable in the model's embedding space?}
			For this question, identifiability is defined by probe accuracy on the detection of typological features of a given language. The input to the probe is the sentence embedding obtained by mean-pooling over the token embeddings of the final layer of the M-BERT model. What constitutes the typological features to be identified may still be subject to change.

			Initially, the probe is trained for language classification. That is, the probing task is a 9-class classification problem where the probe is given a sentence from any of the languages used and it is tasked with identifying the correct language. In this case, each language is viewed as a set of its typological features, and how easily identifiable its features are is quantified by how easily identifiable the language as a whole is. \textbf{Q4} would then be answered in the affirmative when the probe attains a higher accuracy on language $l$ compared to the other languages when applied to $M_l$.

			If time permits, a separate probe might serve to identify specific typological features. The hidden layer would then be followed by separate binary classifiers for each feature in the URIEL \texttt{syntax\_knn} feature vectors. This would more closely mirror the approach taken by Choenni \& Shutova~\cite{choenni-shutova-2022-investigating}. It may well be the case that cross-lingual transfer is not found to be effective between a pair of very similar languages. Should the feature-level probe find that the features shared by that pair are not encoded by the model, that would help explain why transfer is inhibited between the two. That said, analysis of these findings would be more time-consuming than those of the language-identification probe, additional caveats of using \texttt{syntax\_knn} feature vectors notwithstanding. Therefore, feature-level probing will only be attempted after language-level probing has been implemented and sufficient time remains.

	\hfill

% The goal of the project is twofold: Firstly, it aims to show how typology relates to model parameter updates in the fine-tuning stage; secondly, it aims to investigate the effect of typological distance on the success of cross-lingual transfer. Section~\ref{sub:analyses} shows how the setup outlined in Section~\ref{sub:training_setup} serves to achieve these goals. Section~\ref{sub:parameters} lists the parameters used, Table~\ref{tab:data} in Appendix~\ref{app:b} lists the specific datasets used in the experiments.

% 	\subsection{Analyses}
% 	\label{sub:analyses}
% 		\subsubsection{Q1}
% 		In the second stage, four model instances optimized for dependency parsing in German is fine-tuned to do dependency parsing in languages with varying levels of similarity to German. Intuitively, the German-language data in the preceding stage would have prepared it for dependency parsing in the closely related Dutch more than in the unrelated Hungarian. H1 would suggest that the model instance that fine-tunes to Dutch sees less extensive parameter udpates than the one that fine-tunes to Hungarian. In addition, H2 suggests that the updates should primarily be to the earlier layers of the model for Dutch, whereas for Hungarian, they would be across all layers.

% 		The fine-tuning setup includes learned scalar weights $\mathbf{w}_i$ that indicate the importance of each layer $i$. Following Tenney et al.~\cite{tenney-etal-2019-bert}, these can be used to detect where model updates took place for each $\ell$. With $\Delta\mathbf{w}^\ell = \mathbf{w}^\ell - \mathbf{w}^{\textit{deu}}$ being the shift in layer importance after fine-tuning on $\ell$, the expectation is that $\Delta\mathbf{w}^{\textit{nld}}_i$ is greater for layers $i$ that deal with lexical information (i.e., the earlier ones~\cite{tenney-etal-2019-bert}), but lower for layers $i$ that deal with higher-level typological features, since those are generally shared between Dutch and German. On the other hand, the values of $\Delta\mathbf{w}^{\textit{hun}}$ are expected to be more uniform, since Hungarian differs significantly from German both lexically and syntactically.

% 		\subsubsection{Q2}
% 		In the third stage, the model instances for languages $\ell$ are fine-tuned further to each the low-resource languages $\lambda$. Each \m{\ell}{\lambda} is then evaluated on the test set for $\lambda$, with performance measured as the Labeled Attachment Score (LAS). H3 suggests that greater cosine similarity $s(\ell, \lambda)$ should correspond to higher performance on $\lambda$'s test set.

% 	\subsection{Parameters}
% 	\label{sub:parameters}
% 		The model is first fine-tuned to Part A\footnote{The HDT treebank training data is split into various parts. With Part A being larger than the entire English-language treebank used by Choenni et al.~\cite{choenni-etal-2023-cross}, the decision was made to use just that part for training.} of the German training set. Choenni et al.~\cite{choenni-etal-2023-cross} use the English EWT treebank with \~12.5k samples in this stage and train for 60 epochs. To match the number of examples in this setup, training on the German treebank spans 11 epochs. In the second stage, the model instances are fine-tuned over 1\,000 iterations with a batch size of 20 for all languages $\ell$. In the third stage, the models are fine-tuned on 20 examples sampled randomly from each treebank of language $\lambda$. For the Upper Sorbian treebank, these are sampled from the train set. The other three treebanks do not split their data into separate sets; therefore, the sampled examples are removed from the dataset so as to prevent evaluating on them. 

% 		Following Choenni et al.~\cite{choenni-etal-2023-cross}, a cosine-based learning rate scheduler with 10\,\% warm-up and the Adam optimizer are used. For the language model, a learning rate of $1e-04$ is used, and for the classifier, a learning rate of $1e-03$ is used.
